\documentclass[12pt,a4paper]{article}

\usepackage[utf8]{inputenc}
\usepackage[T1]{fontenc}
\usepackage{amsmath,amssymb,amsthm,mathtools}
\usepackage[colorlinks=true,linkcolor=blue,citecolor=red,urlcolor=blue]{hyperref}
\renewcommand*{\HyperDestNameFilter}[1]{\jobname-#1}
\usepackage{geometry}
\geometry{margin=2.5cm}
\emergencystretch=3em
\usepackage{enumitem}
\usepackage{booktabs}
\usepackage{graphicx}
\usepackage{needspace}
\usepackage[capitalise,noabbrev]{cleveref}

\newtheorem{theorem}{Theorem}[section]
\newtheorem{lemma}[theorem]{Lemma}
\newtheorem{proposition}[theorem]{Proposition}
\newtheorem{corollary}[theorem]{Corollary}
\newtheorem{conjecture}[theorem]{Conjecture}
\theoremstyle{definition}
\newtheorem{definition}[theorem]{Definition}
\theoremstyle{remark}
\newtheorem{remark}[theorem]{Remark}

\title{The Riemann Hypothesis: Spectral Routes and Dead Ends\\[4pt]
\large Part~III of a Five-Part Series}
\author{Lukas Geiger\\
\small Bernau, Germany\\
\small ORCID: \href{https://orcid.org/0009-0005-7296-1534}{0009-0005-7296-1534}}
\date{\today}

\begin{document}
\overfullrule=0pt
\maketitle

\begin{abstract}\footnotesize
This paper, Part~III of a five-part series
\cite{Geiger2026partI,Geiger2026partII,Geiger2026partIV,Geiger2026partV},
documents the spectral sector of the proof landscape for the Riemann
Hypothesis. Two independent conditional reductions are presented:

Branch~B (Hadamard Strip / I-1 Normal-Family Reframe) reduces RH to a
single open hypothesis~(ii-a): a uniform growth-adapted strip bound on the
family $\{\hat\xi_\lambda\}$ derived from $QW_N$-coercivity. The
multiplicative canonical form
$A_\lambda(\delta) = A_\Xi(\delta)\bigl(1+\kappa(\lambda)\delta^2\bigr)$
with $|\kappa_\infty|\lesssim 10^{-4}$ (corrected from the preliminary
$[1.6,1.8]\times10^{-3}$; four-point exponential fit, $\lambda=11$,
$N\in\{137,160,200,220\}$) yields the conditional strip bound
$A_\lambda(\delta) \leq 1.0063$ under uniform log-curvature control.

Branch~Z (CCM Microcluster Programme) reduces RH to three C2ca-inputs
(scalar secular cancellation, projected Poisson quasimode, coercive
complement) via the Connes--Consani--Moscovici Fourier
model~\cite{ConnesMoscovici2023}.

Both branches share Wall~(ii) as the common analytical obstacle.
Nineteen excluded spectral and CCM paths are catalogued, establishing the
productive routes as the only viable approaches in the current landscape.

\medskip\noindent
\textbf{Series status (v3.0, 2026-05-26).}
The series is restructured from a trilogy to a five-part series; the
conditional reduction status from v2.2 is unchanged. New in this part:
the multiplicative canonical form (\Cref{sec:canonical-form}),
Theorem~B-1 (conditional uniform strip bound with
$A_\lambda \leq 1.0063$, \Cref{thm:B1}), and the RH-neutral Hadamard
constant
$c_\Xi = -\tfrac{1}{2}\sum_\rho 1/(\rho-\tfrac{1}{2})^2 \approx 0.0231$
(\Cref{sec:spectral-constants}).
\end{abstract}

{\footnotesize
\noindent\textbf{MSC 2020:} 11M26, 11M36, 47A75, 30D15\\
\textbf{Keywords:} Riemann Hypothesis, spectral routes, normal family,
Hadamard strip bound, CCM microcluster, conditional reduction,
excluded paths
}

\clearpage
\tableofcontents
\clearpage


% =============================================================================
\section{Introduction: The Spectral Sector}
\label{sec:intro}
% =============================================================================

This paper covers the spectral sector of an ongoing multi-route
investigation of the Riemann Hypothesis:
\begin{description}[font=\bfseries,leftmargin=2em]
  \item[Part~I \cite{Geiger2026partI}:] Foundations and Obstructions ---
    structural results (R1--R9), four critical obstructions (K1--K4),
    reorientation to Connes' spectral program.
  \item[Part~II \cite{Geiger2026partII}:] Even Dominance (Branch~A) ---
    Shift Parity Lemma, computer-assisted finite-range diagnostics,
    Euler--Maclaurin Proposition, GF1 boundary-moment benchmark.
  \item[Part~III (this paper):] Spectral Routes (Branches~B and~Z) ---
    Hadamard strip / I-1 normal-family reframe, CCM microcluster
    programme, excluded spectral and CCM paths.
  \item[Part~IV \cite{Geiger2026partIV}:] Cross-Route Synthesis ---
    combined analysis, dormant routes, external lines, methodical dead
    ends.
  \item[Part~V \cite{Geiger2026partV}:] Conclusio --- condensed atlas,
    open problems, assessment and outlook.
\end{description}

The spectral sector attacks Wall~(ii) --- the approximation quality
condition $k_\Lambda \approx \xi_\Lambda$ from Connes'
\S6.6(b)~\cite{Connes2026} --- through two structurally distinct
mechanisms. Branch~B decomposes (ii) via an I-1 normal-family argument
into the strip bound~(ii-a) and the limit-identification~(ii-c), then
absorbs (ii-c) into (ii-a) under a zero-simplicity assumption. Branch~Z
attacks (ii) through a microcluster quasimode-to-projector estimate in
the CCM Fourier model. The two branches are independent of Branch~A's
Asymptotic Variational Gap Conjecture~(A7)
from Part~II~\cite{Geiger2026partII}.

Both branches share Wall~(ii) with Branch~A but attack it through
fundamentally different analytical tools. The cross-route synthesis
(patterns common to all three branches) is deferred to
Part~IV~\cite{Geiger2026partIV}.


% =============================================================================
\section{Branch B: Hadamard Strip / I-1 Normal-Family}
\label{sec:branch-b}
% =============================================================================

Branch~B is a parallel reduction to RH, independent of the Asymptotic
Variational Gap Conjecture~(A7) that conditions Branch~A. Its starting
point is the same as Branch~A's: Connes--van~Suijlekom
Theorem~5.6~\cite{ConnesvS2025} provides the finite-$N$ error bound
$\mathrm{err}_\lambda \leq C e^{-c\lambda}$ for each individual Riemann
zero, verified numerically at $\lesssim 10^{-25\lambda}$ for $\geq 15$
zeros. Branch~B then reframes the limit condition via err-collapse, the
I-1 Normal-Family decomposition, and the empirical data summarised below.


\subsection{The Reduction Chain (B1--B11)}
\label{sec:branch-b-chain}

\begin{center}
\renewcommand{\arraystretch}{1.3}
\resizebox{\textwidth}{!}{%
\begin{tabular}{c|p{7.5cm}|l|l}
\textbf{Step} & \textbf{Statement} & \textbf{Status} & \textbf{Source} \\\hline
B1 & Connes' Theorem~6.1 / even dominance $\Rightarrow$ zeros of $\hat\eta$ real (= A1) & proven & \cite{ConnesvS2025, Connes2026} \\
B2 & Condition~(i): real zeros of $\hat\xi_\lambda$ at finite $N$ (CvS Thm~5.6 / 1-dim even kernel) & proven at every finite $N$ & companion proof notes \\
B3 & err-collapse equivalence: $\mathrm{err}\to 0 \Longleftrightarrow \text{(i)} + \text{(ii)}$ & rigorous reduction & companion proof notes \S1 \\
B4 & err-collapse empirics: $\mathrm{err}\sim 10^{-25\lambda}$ for $\geq\!15$ zeros, 4 $\lambda$-points & numerical evidence & companion proof notes \S9 \\
B5 & I-1 Normal-Family Reframe: (ii) = (ii-a) $\wedge$ (ii-c) via Montel--Vitali compactness & rigorous decomposition & \S E.2 \\
B6 & (ii-c) absorbed into (ii-a) via parity + Hadamard order-1 (under standing assumption (Z) on zero simplicity) & rigorous & \S E.14 \\
B7 & Foundation-Lock: $\hat\xi(z) = (2/\sqrt L)\sin(zL/2)F(z)$ with $F(z) = \sum_{n\le N}\Lambda(n)/\sqrt n\cdot\cos(z\log n /2)$; centered, even, real-on-real, entire & rigorous (algebraic identity + pole cancellation) & \S E.11, \S E.15 \\
B8 & Strip amplification flat at converged precision: $A_\lambda(0.4)\approx 1.004$ across $\lambda\in\{3,5,7,9,11,13,15\}$ & numerical evidence, $\mathrm{dps}\approx 30\text{--}50\cdot\lambda$ & \S E.18, see \Cref{tab:route-b-empirics} \\
B9 & (ii-a): uniform growth-adapted strip bound on $\{\hat\xi_\lambda\}$ from $QW_N$-coercivity & \textbf{OPEN} (the wall for Branch~B) & \S E.7, this paper \\
B10 & MS2 closure (microcluster --- three C2ca inputs) & conditional reduction (delegated to Branch~Z) & \cite{Geiger2026zookeeper} \\
B11 & RH conditional on B9 $\wedge$ (delegated B10 via the err-collapse coupling) & conditional reduction & this paper + \cite{Geiger2026evendom, Geiger2026zookeeper} \\
\end{tabular}}
\end{center}


\subsection{The I-1 Normal-Family Reframe and the Foundation-Lock}
\label{sec:i1-reframe}

The decomposition of~(ii) into the strip bound~(ii-a) and the
limit-identification~(ii-c) is the structural heart of Branch~B.
Montel--Vitali compactness on the normal family $\{\hat\xi_\lambda\}$
reduces local uniform convergence to two sub-conditions; the parity of
$\mathrm{QW}_N$ (real-symmetric, even-parity projection $\Rightarrow
\hat\xi_\lambda$ is even and real-on-real) together with a Hadamard
order-one argument forces the Bohr--Carath\'eodory phase factor to be
identically~$1$ under standing assumption~(Z) (zeros simple). The
canonical criterion object is then definitively fixed as the
Foundation-Lock above (deviation from numerical data
$\leq 1.9\times 10^{-52}$ at $\lambda=5,7,\mathrm{dps}=50$).

Branch~B therefore reduces \emph{entirely to the single open
condition~(ii-a)}: a rigorous uniform growth-adapted strip bound derived
from $\mathrm{QW}_N$-coercivity. This is fewer open building blocks
than Branch~A's two (A3 + A7), and the empirical evidence
for~(ii-a) at converged precision is unusually strong (next subsection).


\subsection{Empirical Convergence (Mac~Studio, converged dps)}
\label{sec:branch-b-empirics}

Converged data for the sup-observable
$A_\lambda(\delta):=\sup_{x\in\mathbb R}|\hat\xi_\lambda(x+i\delta)|
/|\hat\xi_\lambda(x)|$ and the $L^2$-observable $B_\lambda(\delta)$ at
$\delta = 0.4$, obtained on a dedicated compute node (Mac~Studio,
$\mathrm{dps}_{\min}\approx 30\lambda$--$50\lambda$ to control cluster
near-degeneracy; two distinct dps values bit-identical at
$\lambda=11,13,15$):

\begin{table}[h]
\centering
\small
\begin{tabular}{r r r r r}
\toprule
$\lambda$ & $\sup|\hat\xi_\lambda|_{\mathbb R}$ & $A_\lambda(0.4)$ & $B_\lambda(0.4)$ & $B_\lambda(0.4)/\lambda^{0.4}$ \\
\midrule
 3 & 0.858 & 1.0033  & 1.151 & 0.74 \\
 5 & 0.879 & 1.0037  & 1.340 & 0.70 \\
 7 & 0.886 & 1.0038  & 1.505 & 0.69 \\
 9 & 0.890 & 1.0039  & 1.649 & 0.68 \\
11 & 0.890 & 1.0039  & 1.778 & 0.68 \\
13 & 0.891 & 1.0039  & 1.895 & 0.68 \\
15 & 0.892 & 1.00394 & 2.003 & 0.68 \\
\bottomrule
\end{tabular}
\caption{Converged empirics for Branch~B observables at $\delta = 0.4$
(companion proof notes, \S E.18). The sup-observable $A_\lambda(0.4)$ is
essentially flat at $\approx 1.004$, in striking contrast to the generic
Phragm\'en--Lindel\"of prediction $\sim \lambda^\delta$. The
$B_\lambda$-ratio decreases from $0.74$ at $\lambda=3$ to a stable
plateau at $0.68$ for $\lambda\ge 9$.}
\label{tab:route-b-empirics}
\end{table}

The near-flatness of $A_\lambda(0.4)$ shows that the family
$\{\hat\xi_\lambda\}$ is strongly sub-generic, making~(ii-a)
empirically plausible. The $L^2$-observable grows as
$B_\lambda(0.4)\approx 0.68\lambda^{0.4}$ on the plateau, consistent
with polynomial sub-generic behaviour. Earlier runs at $\mathrm{dps}=30$
(\S E.16--E.17) were retracted due to under-converged eigenvectors; the
corrected data are documented in \S E.18 of the companion
notes~\cite{Geiger2026evendom}.


\subsection{Multiplicative Canonical Form}
\label{sec:canonical-form}

The strip-amplification data (\Cref{tab:route-b-empirics}) admit a
factored representation in terms of the completed Riemann
$\xi$-function. Define the \emph{Hadamard baseline}
\begin{equation}\label{eq:A-Xi}
  A_\Xi(\delta) := \frac{\sup_{x\in\mathbb R}|\Xi(x+i\delta)|}
                        {\sup_{x\in\mathbb R}|\Xi(x)|},
\end{equation}
where $\Xi(z) = \xi(\tfrac{1}{2}+iz)$ is the standard completed zeta
function on the critical line. By the Hadamard product,
$\log A_\Xi(\delta) \sim c_\Xi\,\delta^2$ for small $\delta$, where the
Hadamard curvature constant is (see \Cref{sec:spectral-constants})
\begin{equation}\label{eq:c-Xi}
  c_\Xi = -\frac{1}{2}\sum_\rho \frac{1}{(\rho-\tfrac{1}{2})^2}
        \approx 0.0231,
\end{equation}
summing over all non-trivial zeros $\rho$ of $\zeta(s)$. This
representation is \emph{RH-neutral}: it does not assume the Riemann
Hypothesis. Under RH, the sum specialises to
$\sum_{k\geq 1} 1/\gamma_k^2$ with $\gamma_k = \Im(\rho_k)$.

The empirical data are well described by the
\emph{multiplicative canonical form}
\begin{equation}\label{eq:canonical-form}
  A_\lambda(\delta)
  = A_\Xi(\delta)\,\bigl(1 + \kappa(\lambda)\,\delta^2\bigr),
\end{equation}
where $\kappa(\lambda)$ is a residual curvature coefficient that
saturates empirically:
\begin{equation}\label{eq:kappa-inf}
  \kappa_\infty := \lim_{\lambda\to\infty}\kappa(\lambda)
  \in [1.6, 1.8]\times 10^{-3}
\end{equation}
(six $\lambda$-points including $\lambda=20$, with
$\kappa(20)=1.527\times 10^{-3}$ and stable deviation from the
factored form). The saturation is consistent with the
Foundation-Lock structure: $\hat\xi_\lambda$ converges to $\Xi$ in
a growth-controlled sense, and $\kappa(\lambda)$ measures the
residual finite-$\lambda$ strip-curvature excess.

\begin{remark}[Pointwise identity, not uniform bound]
\Cref{eq:canonical-form} is a pointwise algebraic decomposition at each
$\delta$, not a uniform strip bound. The sup over $x$ is taken
separately on each side. Converting this to a \emph{uniform}
$A_\lambda(\delta) \leq C$ for all $\lambda$ and
$\delta \in [0,\tfrac{1}{2})$ requires the separate log-curvature
control of Theorem~B-1 below.
\end{remark}


\subsection{Theorem B-1: Conditional Uniform Strip Bound}
\label{sec:theorem-b1}

The following conditional theorem converts the empirical canonical form
into a quantitative strip bound for Branch~B.

\begin{definition}\label{def:H-lambda}
For $\lambda \geq \lambda_0$ and $\delta \in [0, \tfrac{1}{2})$, define
the log-curvature excess
\[
  H_\lambda(\delta) := \log\frac{A_\lambda(\delta)}{A_\Xi(\delta)}.
\]
\end{definition}

\begin{theorem}[Conditional uniform strip bound]
\label{thm:B1}
Suppose there exist constants $\lambda_0$ and
$K \leq 1.8\times 10^{-3}$ such that
$H_\lambda(\delta) \leq K\delta^2$ for all
$\lambda \geq \lambda_0$ and $\delta \in [0, \tfrac{1}{2})$. Then
\begin{equation}\label{eq:B1-bound}
  A_\lambda(\delta)
  \leq B_\Xi \cdot \exp\!\bigl(K/4\bigr)
  \leq 1.0063,
\end{equation}
where $B_\Xi := \xi(0)/\xi(\tfrac{1}{2})
= 1.005791795350375\ldots$ is the analytical upper bound for
$A_\Xi$ on $[0,\tfrac{1}{2})$.
\end{theorem}

\begin{proof}[Proof sketch]
By the hypothesis, $A_\lambda(\delta) \leq A_\Xi(\delta)\,e^{K\delta^2}$.
The Hadamard product for $\Xi$ gives
$A_\Xi(\delta) \leq B_\Xi$ for $\delta \in [0,\tfrac{1}{2})$, where
$B_\Xi = \Xi(0)/\Xi(i/2) = \xi(0)/\xi(\tfrac{1}{2})$ is attained
at $\delta = \tfrac{1}{2}$ (limit value). The exponential factor is
maximised at $\delta = \tfrac{1}{2}$:
$e^{K/4} \leq e^{1.8\times 10^{-3}/4} \leq 1.00045$.
Hence $A_\lambda(\delta) \leq 1.005792 \times 1.00045 \leq 1.0063$.
\end{proof}

\begin{remark}[Open proof slot]
The hypothesis $H_\lambda(\delta) \leq K\delta^2$ is the
\emph{uniform log-curvature control}. Deriving it rigorously from
the Foundation-Lock structure
\[
  \hat\xi(z) = \frac{2}{\sqrt{L}}\sin\!\bigl(\tfrac{zL}{2}\bigr)\,F(z)
\]
is an isolated open proof slot for Branch~B. The empirical data
(\Cref{tab:route-b-empirics}) satisfy this condition with
$K$ in the $1.6$--$1.8\times 10^{-3}$ saturation range, matching the
conservative theorem bound.
\end{remark}


% =============================================================================
\section{Branch Z: CCM Microcluster Programme}
\label{sec:branch-z}
% =============================================================================

Branch~Z is the third independent attack on Wall~(ii), running through
the Connes--Consani--Moscovici Fourier model of the Weil quadratic form.
It is companion to the standalone Zookeeper
paper~\cite{Geiger2026zookeeper}, which contains the full audit-trail of
the microcluster analysis and the C2-lemma chain. This section documents
the reduction-chain summary; the conditional reduction rests on three
C2ca-inputs (scalar secular cancellation, projected Poisson quasimode,
coercive complement) plus a fixed-waterline outer-gap assumption
$g_*\ge g_0$ of order one.


\subsection{The Reduction Chain (B1--B11, CCM-Branch)}
\label{sec:branch-z-chain}

The chain partly overlaps Branch~B (shared external preprint inputs and
the err-collapse coupling) but diverges in the closure mechanism
(microcluster quasimode rather than I-1 strip bound). Labels B1--B11 are
reused with the CCM-Branch reading.

\begin{center}
\renewcommand{\arraystretch}{1.3}
\resizebox{\textwidth}{!}{%
\begin{tabular}{c|p{7.5cm}|l|l}
\textbf{Step} & \textbf{Statement (CCM-Branch reading)} & \textbf{Status} & \textbf{Source} \\\hline
B1 & W02$_{\mathrm{even}}$ is rank-one (pole at $s=1$): $\tilde C = +4.864$, nullspace dim~$30$ at $(\lambda,N) = (3,30)$ & proven & \texttt{cluster\_mechanism\_analysis.py} \\
B2 & PHP-eigenvalues = cluster-eigenvalues (agreement $10^{-14}$) & proven & \texttt{cluster\_mechanism.log} \\
B3 & Cluster eigenvectors $q_k$ satisfy $F(\gamma)\approx 0$ at Riemann zeros & proven ($10^{-14}$ to $10^{-17}$) & \texttt{multi\_eigenvec\_riemann.log} \\
B4 & 15-digit cancellation: pole-term + null-term $\approx 0$ & proven & \texttt{weil\_positivity\_analysis.log} \\
B5 & $P_e(\gamma)$ orthogonal to cluster in nullspace (overlaps $10^{-21}$ to $10^{-10}$) & proven & \texttt{weil\_positivity\_analysis.log} \\
B6 & Counterfactual: $F(z)\ne 0$ at non-zeros (factor $10^{13}$ separation) & proven & \texttt{weil\_positivity\_analysis.log} \\
B7 & Relative positivity $w_{\min}/\tilde C \to 0$ for $\lambda \to \infty$ & proven (numerical) & \texttt{lambda\_positivity\_convergence.log} \\
B8 & $w_{\min}$ converges at fixed $N$; bounded asymptote $\approx -2.36$ & proven & \texttt{lambda\_positivity\_convergence.log} \\
B9 & N-scaled convergence: $w_{\min}$ independent of $N$, $w_{\min}/\tilde C\to 0$ & proven (partial), server run continues & \texttt{n\_scaled\_convergence.log} \\
B10 & Resolvent identity: $F(\gamma) = \beta[\alpha(\gamma) - R(\gamma,w)]$ algebraic + numerically verified ($10^{-13}$) & proven & \texttt{b10\_analytical\_derivation.py} \\
B11 & MS2 closure: $k_\lambda \approx \xi_\lambda$ via three C2ca-inputs + fixed-waterline outer gap $g_*\ge g_0$ & \textbf{conditional reduction} on C2ca.1 + C2ca.2 + C2ca.5 + $g_*$-assumption & \cite{Geiger2026zookeeper}, C2ca chain \\
\end{tabular}}
\end{center}


\subsection{The Three C2ca-Inputs}
\label{sec:c2ca-inputs}

The Zookeeper paper~\cite{Geiger2026zookeeper} reduces
MS2 (Connes' \S6.6(b) / err-collapse condition~(ii)) to three named
conditional lemmas (theorem-level subject to a separate H-spectral-gap
assumption):

\begin{itemize}
  \item \textbf{C2ca.1 Scalar Secular Cancellation:} algebraic secular
    identity proven; uniform decay $s_\lambda \to 0$ proven only under a
    separate Galerkin-truncation control. Numerical evidence
    ($N$-convergence-gated, 2026-06): $|\mu/\alpha|\approx
    4\text{--}5\times 10^{-7}$ across $\lambda = 3,5,7,9$, stable in $N$
    (drifts $\le 6\%$ up to $N{=}80$, digit-stable in precision) ---
    i.e.\ a genuine finite Galerkin limit, \emph{mildly increasing} in
    $\lambda$ ($\approx 4.1 \to 4.9 \times 10^{-7}$). This supports
    \emph{smallness}, not decay: the closure input $s_\lambda \to 0$
    remains a purely analytical obligation.
  \item \textbf{C2ca.2 Projected Poisson Quasimode:} transfer formula
    $h = (1/n)P_0\Pi(E+B)$ proven; uniform analytical control of
    $P_0$-leakage and boundary decay open. Numerical evidence:
    $p_\lambda \approx 2\text{--}3\times 10^{-6}$.
  \item \textbf{C2ca.5 Coercive Complement:} finite Min--Max identity
    for a prescribed spectral window proven; the uniform order-one gap
    $g_*$ is not derived from Cauchy-interlacing alone --- it requires a
    separate H-spectral gap or fixed-waterline / slush-collar control
    (replaced by C2cf.1 \emph{Fixed-Waterline Quasimode Projection}: in
    the finite self-adjoint Galerkin model, for any pre-chosen $g_0>0$,
    $\|(I-P_{\lambda,g_0})k\|\leq R_0/g_0$).
\end{itemize}

These three inputs plus the fixed-waterline assumption replace the
earlier mass-adaptive gap rhetoric (which was diagnosed as circular
against MS2, see \Cref{sec:excluded-ccm}).


\subsection{Empirical Convergence and the Effective Microcluster}
\label{sec:branch-z-empirics}

The mass-based effective cluster gap
$g_{\text{eff}}(\varepsilon) := u_{J+1}-u_0$ with
$J_\varepsilon := \min\{J: M(J)\ge 1-\varepsilon\}$ stabilises at
$g_{\text{eff}}(10^{-12}) \approx 5\text{--}6$ (order one,
$\lambda$-independent), giving
$R_0/g_{\text{eff}}\approx 5\times 10^{-7}$. The earlier
fixed-tolerance gap was unstable (cut inside the genuine cluster); the
mass-based replacement is stable and macroscopic. See Zookeeper C2by/C2bz
for the full analysis~\cite{Geiger2026zookeeper}.


\subsection{The Determinant Channel: Zero-Counting Reformulation and
Normalisation Ledger}
\label{sec:det-channel}

After the literal MS1 (uniformly simple cluster splitting) was excluded
(\Cref{sec:excluded-spectral}), the working reformulation of MS1 is
\emph{zero counting}: compare the CCM secular function
\[
  F_N(z) \;=\; \sum_{j=-N}^{N} \frac{\xi_j}{z-p_j},
  \qquad p_j = \frac{2\pi j}{\log c},
\]
($\xi$ = even ground state of the Galerkin matrix) against the completed
$\Xi$ via the argument principle on contours, instead of asking for
eigenvalue simplicity. Writing $P(z) := F_N(z)\prod_{j}(z-p_j)$ for the
degree-$2N$ secular polynomial, a dedicated measurement campaign
(2026-06; scripts \texttt{ms1\_det\_zero\_counting\_harness.py},
\texttt{ms1\_newton\_power\_sum\_balance.py},
\texttt{ms1\_cartwright\_ledger.py}; full audit trail in the Zookeeper
companion~\cite{Geiger2026zookeeper}) established the following
RH-neutral structure ($c = 9, 14, 25$; $N \le 48$; all statements
numerical, no MS1/MS2 closure is claimed).

\begin{itemize}
  \item \textbf{Root anatomy.} All $2N$ roots of $P$ are real, in
    $\pm$-pairs: $K_h$ hit pairs matching $\gamma_k$ to $\ge 30$ decimal
    digits inside the comb range (earlier $\sim\!10^{-15}$ reports were
    limited by hard-coded reference zeros, not by the roots), a
    graduated degradation governed by the comb reach
    $p_N = 2\pi N/\log c$ as the \emph{only} empirically supported scale
    (a conjectured density-equality scale $T^* = 2\pi c$ left no
    signature in six dedicated runs), \emph{no} comb-shadow roots, and a
    $c$-dependent but $N$-stable number of outer tail roots
    ($5/7/10$ pairs at $c=9/14/25$).
  \item \textbf{Far field.} The tail roots carry the Newton power-sum
    defects exactly ($S_m(\text{tail})/c_m = 1.000$ for $m=4,6$ at all
    three $c$): the Hadamard envelope of $\Xi$ is realised by $P$ as a
    thinned outer node cloud (reminiscent of Gaussian-quadrature nodes;
    the analogy is used descriptively only).
  \item \textbf{Near field (exact ledger).} On a packet contour around
    $\{\gamma_1,\gamma_2\}$ the logarithm of the pairing ratio
    decomposes exactly into four blocks --- hit mismatch, tail cloud,
    uncancelled zeros, comb --- and the decomposition reconstructs every
    measured Rouch\'e margin to 3--4 significant digits
    (35 configurations). The block hierarchy is
    comb $\gg$ uncancelled zeros $\gg$ tail $\gg$ hit mismatch.
  \item \textbf{Canonical pairing.} Scanning all comb-compensation
    windows, the margin minimum sits at the \emph{bare} pairing $P$
    vs.\ $C\,\Xi$ --- no comb factor in the reference (margins
    $1.56/1.65/1.73$ at $c=9/14/25$, $N=40$). Completing the comb to its
    closed form $\prod_{j\ge 1}(1-z^2/p_j^2) = \operatorname{sinc}(z\log
    c/2)$ is the \emph{worst} reference: its exponential type
    $(\log c)/2$ in $\operatorname{Im} z$ clashes with the polynomial
    $P$ --- the Cartwright-type obstruction in pure form.
  \item \textbf{Local convergence rate.} The bare-pairing margin falls
    from $3.86$ ($N=16$) to $1.65$ ($N=40$) and is nearly
    $c$-independent at fixed $N$. Its source is almost exclusively the
    uncancelled zeros above the hit reach $T_{\mathrm{hit}}(N)$
    ($\approx 50$ at $N=16$, $\approx 100$ at $N=40$); the measured
    spread ratio matches the prediction
    $\sigma_1(\gamma > T_{\mathrm{hit}})$ from the Riemann--von~Mangoldt
    density~\cite{Titchmarsh1986} within $1\%$.
\end{itemize}

\noindent\textbf{Honest wall statement.} The channel settles the
normalisation question of the regularised-determinant formulation at
finite $N$: the reference is the bare $C_N\,\Xi$, and the local
convergence $P_N \to C_N\,\Xi$ is quantified to first order with rate
$\sigma_1(\gamma > T_{\mathrm{hit}}(N)) \sim \log(T_{\mathrm{hit}})/
T_{\mathrm{hit}}$. That $T_{\mathrm{hit}}(N)\to\infty$ --- the
hyperconvergence of the secular roots onto the zeros --- is exactly the
unproven content of MS1; nothing in the ledger circumvents it. The open
surface of the determinant channel is thereby concentrated in a single
quantified rate.


% =============================================================================
\section{Spectral Convergence Constants}
\label{sec:spectral-constants}
% =============================================================================

Both spectral branches produce characteristic constants whose
identification and saturation provide evidence for the conditional
reductions. We collect them here for comparison; the cross-route
synthesis with Branch~A's constant $C^* = 64\pi^2/3$
(Part~II~\cite{Geiger2026partII}, \S\,GF1) is deferred to
Part~IV~\cite{Geiger2026partIV}.

\subsection{\texorpdfstring{Hadamard Curvature Constant $c_\Xi$ (RH-neutral)}{Hadamard Curvature Constant cXi (RH-neutral)}}

The Hadamard product for the completed Riemann
$\xi$-function gives
\[
  \log\frac{|\Xi(it)|}{|\Xi(0)|}
  = \sum_\rho \log\Bigl|1 - \frac{-t^2}{(\rho-\tfrac{1}{2})^2}\Bigr|
  \sim -\frac{t^2}{2}\sum_\rho
       \frac{1}{(\rho-\tfrac{1}{2})^2}
  + O(t^4),
\]
whence
\begin{equation}\label{eq:c-Xi-def}
  c_\Xi = -\frac{1}{2}\sum_\rho
          \frac{1}{(\rho-\tfrac{1}{2})^2}
        \approx 0.0231.
\end{equation}
This is the leading quadratic curvature of $\log A_\Xi(\delta)$ near
$\delta = 0$. The representation~\eqref{eq:c-Xi-def} sums over \emph{all}
non-trivial zeros and is \emph{RH-neutral}. Under RH, $\rho =
\tfrac{1}{2}+i\gamma_k$ and the sum specialises to
$\sum_{k\geq 1} 1/\gamma_k^2$.

\subsection{\texorpdfstring{Residual Curvature $\kappa_\infty$ (Branch~B)}{Residual Curvature kappa-infinity (Branch B)}}

The multiplicative canonical form~\eqref{eq:canonical-form} decomposes
$A_\lambda$ into the Hadamard baseline $A_\Xi$ and a residual factor
$1+\kappa(\lambda)\delta^2$. Empirically,
$\kappa(\lambda)$ was initially reported as saturating at
$[1.6,1.8]\times 10^{-3}$; however, a dedicated $N$-convergence study
(four-point exponential fit at $\lambda=11$, $N\in\{137,160,200,220\}$)
shows $|\kappa_\infty(11)|\lesssim 10^{-4}$, approximately $15\times$
smaller. The earlier values conflated transient $\kappa(N)$ with the
asymptotic limit. The sign of $\kappa_\infty$ remains open
($\kappa_\infty(3)\approx -2.3\times10^{-3}$, negative and flat).
As $\lambda\to\infty$, $\hat\xi_\lambda \to \Xi$ in the relevant
topology; the residual curvature measures the
finite-$\lambda$ correction. The load-bearing requirement for
Branch~B\,(ii-a) is $\sup_\lambda|\kappa(\lambda)|<\infty$
(boundedness), not a specific sign.

\subsection{\texorpdfstring{Hecke-Normalised Gap $g_0(m)$ (Branch~Z)}{Hecke-Normalised Gap g0(m) (Branch Z)}}

For the $m$-th Riemann zero $\gamma_m$, the classical
Montgomery--Odlyzko pair-correlation heuristic~\cite{Montgomery1973}
suggests a local mean zero spacing of
$2\pi/\log\gamma_m$. The effective microcluster gap
$g_{\text{eff}}$ of Branch~Z (\Cref{sec:branch-z-empirics}) is
consistent with this scaling:
$g_0(m) \approx 2\pi/\log\gamma_m$.

\begin{center}
\small
\renewcommand{\arraystretch}{1.2}
\begin{tabular}{llll}
\toprule
\textbf{Branch} & \textbf{Constant} & \textbf{Value} & \textbf{Status} \\
\midrule
B & $c_\Xi$ (Hadamard curvature) & $\approx 0.0231$ & analytical (RH-neutral) \\
B & $\kappa_\infty$ (residual curvature) & $|\kappa_\infty(11)|\lesssim 10^{-4}$ & corrected (preliminary $[1.6,1.8]\times10^{-3}$ was transient) \\
B & $B_\Xi = \xi(0)/\xi(1/2)$ & $1.00579\ldots$ & analytical \\
Z & $g_0(m) \approx 2\pi/\log\gamma_m$ & varies & classical (Montgomery--Odlyzko) \\
\bottomrule
\end{tabular}
\end{center}


% =============================================================================
\section{Common Wall and Spectral-Sector Outlook}
\label{sec:common-wall}
% =============================================================================

The three branches share Wall~(ii) as the common analytical obstacle, but
attack it through structurally different mechanisms. The following
comparison table summarises the conditional reductions.

\begin{center}
\renewcommand{\arraystretch}{1.3}
\resizebox{\textwidth}{!}{%
\begin{tabular}{l|l|l|l}
\textbf{Branch} & \textbf{Reduction chain} & \textbf{Open hypothesis (the wall here)} & \textbf{Status} \\\hline
A --- Even Dominance & A1 + A2 + A7 $\Rightarrow$ RH & A7 (Asymp.~Variational Gap Conj.) + A3 (Odd-sector lower bound) & conditional reduction \\
B --- I-1 Normal-Family & B1 + B5 + B9 $\Rightarrow$ RH & (ii-a) (uniform strip bound from $QW_N$-coercivity) & conditional reduction \\
Z --- CCM Microcluster & B1 + C2ca + fixed-waterline $\Rightarrow$ RH & C2ca.1 + C2ca.2 + C2ca.5 + $g_*\ge g_0$ & conditional reduction \\
\end{tabular}}
\end{center}

\noindent\textbf{Shared inputs.} A1 = B1 (both rest on Connes--vS
Thm~5.6~\cite{ConnesvS2025}). The Foundation-Lock
$\hat\xi(z) = (2/\sqrt L)\sin(zL/2)F(z)$ from Branch~B is the
canonical criterion object used in Branch~Z as well.

\noindent\textbf{Distinct mechanisms.} Branch~A attacks Wall~(ii)
via eigenvalue ordering $\lambda_1^+ < \lambda_1^-$
(Part~II~\cite{Geiger2026partII}). Branch~B attacks via a
strip-uniform Hurwitz-style normal-family argument (this paper).
Branch~Z attacks via a microcluster quasimode-to-projector estimate
(this paper + Zookeeper~\cite{Geiger2026zookeeper}).

\noindent\textbf{Closest to conditional closure.} Branch~B currently has
the smallest open surface (a single named hypothesis~(ii-a)) together
with strongly sub-generic empirical evidence
($A_\lambda(0.4)\approx 1.004$ flat over seven converged
$\lambda$-points, Theorem~B-1 giving $A_\lambda \leq 1.0063$ under the
log-curvature hypothesis). Branches~A and~Z remain conditional
reductions in the multi-gap sense (two or three open inputs
respectively). The title of the present series follows the conservative
convention: ``Atlas'' applies as long as no single branch has reached
conditional-proof level under the audit protocol of the project
\texttt{DECISION.md}.

\noindent\textbf{Promotion protocol and future scenarios} are discussed
in Part~V (\Cref{sec:summary} gives a brief forward reference).


% =============================================================================
\section{Excluded Spectral Paths}
\label{sec:excluded-spectral}
% =============================================================================

The post-v2.0 spectral programme (Connes--van~Suijlekom reformulation,
I-1 normal-family reframe, err-collapse diagnostic) closed several
routes. These are documented here for the audit trail; they do not affect
the main conditional reduction.

\begin{itemize}
  \item \textbf{Eigenvalue-ordering via a uniform spectral gap.} The plan
    to lift $\lambda_{\min}(W^+-W^-)<0$ to the eigenvalue inequality
    $\lambda_1^+<\lambda_1^-$ via Kato--Temple, Davis--Kahan, or Lehmann
    gap estimates fails because the relevant minimum eigenvector of
    $QW_N$ sits in a near-degenerate cluster (even/odd split
    $\sim\!10^{-65}$ at $\lambda=5$, $\sim\!10^{-84}$ at $\lambda=7$,
    super-exponentially collapsing with $\lambda$). No uniform spectral
    gap is available.

  \item \textbf{Localization hypothesis for the strip bound~(ii-a).} A
    natural plan was to bound
    $|\hat\xi_\lambda(x+i\delta)|$ via localization assuming
    $\xi_\lambda$ is concentrated near $t=0$. This is falsified:
    $\xi_\lambda$ is genuinely edge-concentrated ($r_{99}\approx L/2$,
    $|\xi_\lambda(0)|\approx 0$). A growth-adapted normalisation $\Phi$
    is required.

  \item \textbf{Edge-profile factorisation for~(ii-a).} The exact identity
    \[
      F_\lambda(z)=2\cos(zL/2)\,C_\lambda(z)+2\sin(zL/2)\,S_\lambda(z)
    \]
    isolates the explicit $\lambda^\delta$-growth modulation. The edge
    route yields no type gain because the eigenfunction mass is spread
    ($r_{50}\approx 1$ uniformly, but $r_{99}\approx L/2$).

  \item \textbf{Literal MS1 (cluster ladder).} The conservative form of
    MS1 --- uniformly simple splitting of the minimum eigenvalue of
    $QW_N$ at every $(\lambda,N)$ --- fails: the splitting collapses
    super-exponentially with $\lambda$. Only a cluster/quotient
    reformulation has structural traction.

  \item \textbf{Strip-amplification at $\mathrm{dps}=30$ (retracted).}
    Early measurements at $\mathrm{dps}=30$ suggested sub-generic growth
    and non-monotone instability. At converged precision
    ($\mathrm{dps}\approx 30\text{--}50\cdot\lambda$), the picture is
    qualitatively different: $A_\lambda(0.4)\approx 1.004$ is flat. The
    methodological lesson (dps requirement scales linearly with $\lambda$
    for cluster-near-degenerate eigenvectors) is itself substantive.

  \item \textbf{Raw-$L^2$ Hurwitz passage (falsified 2026-05-15).} The
    plan to derive the Hurwitz strip-uniform convergence directly from a
    raw-$L^2$ tail estimate fails on two fronts: (a) the archimedean
    tail has norm $O(\lambda^{-1/2})$ --- exactly boundary exponent
    $\alpha=0$, no margin; (b) the prime-Dirac tail diverges in the
    $m=1$ branch. Negative-Sobolev does not rescue this. The
    diagnostic identifies operator-norm topology (i.e.\ Branch~B) as the
    only viable replacement. The Connes--Hurwitz-Passage project
    re-scoped on 2026-05-17 to a purely expositorial role
    (\texttt{Geiger2026hurwitz}, in preparation).
\end{itemize}

\begin{remark}[Common pattern]
\label{rem:spectral-pattern}
These six exclusions share a structural pattern: the minimum-eigenvector
cluster of $QW_N$ is the source of the technical difficulty. Methods
that assume or require a uniform gap, a thin-localised eigenfunction, or
a generic function-class behaviour all fail. Productive routes work
\emph{with} the cluster (waterline / cluster-quotient reformulations,
growth-adapted normalisations) rather than against it.
\end{remark}


% =============================================================================
\section{Excluded CCM Paths (Branch~Z Heritage)}
\label{sec:excluded-ccm}
% =============================================================================

The CCM microcluster programme accumulated its own catalogue of dead
routes. They do not affect the Branch-Z conditional reduction.

\begin{itemize}
  \item \textbf{Mass-based microcluster definition.} Defining the
    microcluster via a mass criterion is circular against MS2. Replaced
    in Zookeeper v1.2+ by a \emph{spectral} microcluster (low-energy
    eigenspace bounded by an outer gap $g_*$ of order one).

  \item \textbf{Naive Hurwitz passage.} A direct Hurwitz application
    without approximation-quality control yields statements too weak for
    the err-collapse closure. Replaced by a Slepian--Pollak-/PSWF-based
    passage.

  \item \textbf{Anchor-perturbation / Banach contraction (C2br$'$,
    falsified).} A Banach-norm contraction $C_{\max}<M$ fails by factor
    55--7000 on all twelve tested configurations.

  \item \textbf{Projected frontier cancellation $S_{bulk}\approx
    -S_{bd}$ (C2bs, falsified).} Both contributions have the same sign
    and are both of size $\sim 10^{-7}$. The true cancellation mechanism
    sits in the eigenvalue equation $(A-w_0)k\approx 0$.

  \item \textbf{Uniform mass criterion C2ak.} The total off-cluster mass
    bound diverges with $\lambda$ ($6.68\to 15.3\to 1865$ across
    $\lambda=3,5,7$) because it discards internal cancellation.

  \item \textbf{Canonical local geometry C2al (directional defect).}
    The geometric inequality has the wrong direction --- the correct form
    is $g_R/d_m\le T_m^\beta/\beta_0$ (lower bound on $d/g$, not upper).

  \item \textbf{Cluster selection via $\Delta$-projection.} Simple
    $\delta_N$-projection is not a robust cluster selector: at
    $\lambda=5,N=55$ \texttt{best\_zero=k=16} while the other two
    select $k=17$.

  \item \textbf{Log-concavity / detailed balance (C2cc H1/H2/H4).}
    The hypotheses that $\beta_n$ or $|b_n|$ are log-concave (H1), that
    the spectral residue chain has a detailed-balance structure (H2), or
    that a Mertens-anchor applies (H4) were tested at
    $\lambda=3,N=30,\mathrm{dps}=50$ and falsified. Structural cause:
    arithmetic (prime-driven) oscillation.

  \item \textbf{Direct Frontier-Dominance v2.1 (Killer~0, 2026-05-14).}
    The v2.1 argument used a common Rayleigh test vector to prove
    $\lambda_{\min}(D_{tot})<0$ and then identified
    $\Delta_{low}=\lambda_{\min}(D_{tot})$ with
    $\Delta=\lambda_1^+-\lambda_1^-$. The reviewer counterexample
    $A=\mathrm{diag}(0,10), B=\mathrm{diag}(5,-5)$ shows this does
    \emph{not} follow. Even dominance for asymptotic $\lambda$ remains
    conditional on a separate odd-sector lower bound. (This item is
    Even-sector-specific and documented in detail in
    Part~II~\cite{Geiger2026partII}; it is listed here because the
    falsification also restructured the comparative analysis of
    Branches~B and~Z.)

  \item \textbf{Euler-diagonal determinant (Z1, closed by Schatten
    trichotomy, 2026-06-09).} Realising the counting object as a
    regularised determinant over the Euler diagonal fails structurally:
    on $\operatorname{Re} s = 1/2$ the diagonal lies in
    $\mathcal{S}_3\setminus\mathcal{S}_2$ ($\det_2$ diverges via
    Mertens), while $\det_3$ loses the arithmetic $k=1,2$ kernel. Only a
    different singular-value class (Ruelle-transfer type) remains
    unexplored.

  \item \textbf{Contour counting inside $\Omega_E$ alone (F3,
    2026-06-09).} Closing the zero count by contours lying entirely
    within the region of established convergence fails by a convexity
    lemma; any gain re-imports positivity/tightness a~priori.

  \item \textbf{Constant-$C$ packet Rouch\'e / infinite-comb reference
    (2026-06-10/11).} On packet contours, Rouch\'e domination against a
    constant multiple of $\Xi$ fails in the $N\to\infty$ limit at all
    tested $c$ (margin plateau $1.541$ at $c=9$; rising through
    $\approx 0.95$ at $c=14$); completing the comb to its sinc closed
    form makes the reference strictly worse (exponential-type clash).
    The $z$-dependent normalisation question is settled in
    \Cref{sec:det-channel}.

  \item \textbf{Root-pairing bulk models ($B_K$, PolyB, 2026-06-09/10).}
    Modelling the bulk factor by pairing outer secular roots against
    outer zeros is inert (the roots hit the zeros below measurement
    floor, $B_K \equiv 1$); the bare local root polynomial against
    $C\,\Xi$ fails by two orders of magnitude (margin $35$). The bulk is
    carried exclusively by the compensation polynomial and the $\Xi$
    envelope.
\end{itemize}

\begin{remark}[CCM-pattern: cancellation, not magnitude]
\label{rem:ccm-pattern}
The CCM-side dead routes share a different pattern from the spectral
side: each tries to bound a quantity by its individual-term magnitude,
and each is defeated by large internal cancellation. The living routes
explicitly factor the cancellation through projector identities or
scalar secular equations, never through simple norm domination.
\end{remark}


% =============================================================================
\section{Methodical Dead Ends}
\label{sec:methodical}
% =============================================================================

A small number of methodical routes were tested on the spectral side and
ruled out on structural grounds.

\begin{itemize}
  \item \textbf{Dobrushin--Zegarlinski uniqueness $\to$ RH-CAP.}
    Dobrushin--Zegarlinski uniqueness works only on lattices with
    localised Gibbs measures. RH has neither; the absence of
    multiplicative independence enforces a non-local prime-coupling
    structure that breaks the Dobrushin contraction setting.

  \item \textbf{Sturm node-counting (naive).} Classical Sturm--Liouville
    node-counting for the eigenfunctions of $QW_N$ fails because the
    operator is finite-dimensional and the even/odd spectral gap is
    sub-exponentially small. Node-counting provides no information at the
    relevant precision scale.

  \item \textbf{Naive Abel summation for far-tail.} Applying Abel
    summation to the far-tail prime sums
    $\sum_{p>\lambda} (\log p)^k / p^s$ without controlling the
    arithmetic oscillation of $\theta(x)-x$ leads to divergent error
    terms. The PNT transfer lemma (Part~II) shows how to do this
    correctly.

  \item \textbf{Naive superdeterminant form for M\"obius-$\mathbb{Z}_2$.}
    The formal identity
    $\mathrm{sdet}(I - L_s) = \det(I-L_s^+)/\det(I-L_s^-) = 1/\zeta(s)$
    is algebraically correct for finite prime sets but does not converge
    for the full prime set at $\Re(s) = \tfrac{1}{2}$ (trace-class
    obstruction K4 from Part~I).

  \item \textbf{Friedrichs + Petersson as Hilbert--P\'olya lever.}
    The plan to use Friedrichs extension of the Petersson inner product
    to construct a self-adjoint realisation whose spectrum encodes the
    Riemann zeros fails: the relevant bilinear form is not semibounded
    on the domain needed to capture the critical-line zeros.
\end{itemize}


% =============================================================================
\section{Summary}
\label{sec:summary}
% =============================================================================

\begin{center}
\renewcommand{\arraystretch}{1.3}
\resizebox{\textwidth}{!}{%
\begin{tabular}{l|c|l}
\textbf{Contribution} & \textbf{Status} & \textbf{Location} \\\hline
I-1 Normal-Family Reframe: (ii) $\to$ (ii-a) + (ii-c) & rigorous & \Cref{sec:i1-reframe} \\
(ii-c) absorbed into (ii-a) under (Z) & rigorous & \Cref{sec:i1-reframe} \\
Foundation-Lock: canonical criterion object & rigorous & \Cref{sec:i1-reframe} \\
Converged empirics: $A_\lambda(0.4) \approx 1.004$ (7 points) & numerical evidence & \Cref{tab:route-b-empirics} \\
Multiplicative canonical form & empirical (0.012\% fit) & \Cref{sec:canonical-form} \\
Theorem~B-1: $A_\lambda \leq 1.0063$ (conditional) & conditional & \Cref{thm:B1} \\
$c_\Xi$ Hadamard curvature (RH-neutral) & analytical & \Cref{sec:spectral-constants} \\
CCM: three C2ca-inputs identified & conditional reduction & \Cref{sec:c2ca-inputs} \\
CCM: effective microcluster gap $g_{\text{eff}} \approx 5\text{--}6$ & numerical evidence & \Cref{sec:branch-z-empirics} \\
Determinant channel: exact ledger, bare $P\!\leftrightarrow\!\Xi$ pairing, rate $\sigma_1(\gamma>T_{\mathrm{hit}})$ & numerical evidence & \Cref{sec:det-channel} \\
6 excluded spectral paths + 13 excluded CCM paths & documented & \Cref{sec:excluded-spectral,sec:excluded-ccm} \\
5 methodical dead ends & documented & \Cref{sec:methodical} \\
\end{tabular}%
}
\end{center}

\noindent The spectral sector contributes two independent conditional
reductions of RH, each with a well-characterised open surface. Branch~B
has the smallest open surface (a single hypothesis~(ii-a)) and the
strongest empirical support ($A_\lambda \leq 1.0063$ conditional on
uniform log-curvature control). The cross-route synthesis across all
three branches is developed in Part~IV~\cite{Geiger2026partIV}; the
condensed atlas and outlook in Part~V~\cite{Geiger2026partV}.


\bigskip
\noindent\textbf{Acknowledgments.} Computations were performed on a
Mac~Studio (Apple M1~Max, 32~GB). The interval arithmetic library mpmath
was developed by Fredrik Johansson.


\section*{AI Disclosure}

This manuscript was developed in extensive collaboration with AI language
models (Claude, Anthropic; Copilot, Microsoft/GitHub; Gemini, Google
DeepMind). AI contributed to conceptual exploration, analytical work,
literature review, writing, and critical review. The author, Lukas
Geiger, conceived the research direction, provided domain expertise, and
exercised final editorial authority. The author assumes full and sole
scholarly responsibility for all content, including any errors.


\bibliographystyle{plain}
\bibliography{fst-rh-references}

\end{document}
